%!TEX root = TheoretischePhysikA.tex

\section{Vorlesung: Lineare Koordinatentransformationen}
Wir betrachten zwei Koordinatensysteme $\Sigma$ und $\Sigma'$, deren Ursprung übereinstimmt und deren Achsen gegeneinander verdreht sind.
\figbox
Der Punkt $P_0$ mit den Koordinaten
\begin{equation}
  \begin{pmatrix}
    1 \\ 0
  \end{pmatrix}
\end{equation}
in $\Sigma'$ hat in $\Sigma$ die Koordinaten
\begin{equation}
  \begin{pmatrix}
    \cos(\theta) \\ \sin(\theta)
  \end{pmatrix}.
\end{equation}
Analog, der Punkt $P_1$ mit Koordinaten in $\Sigma'$
\begin{equation}
  \begin{pmatrix}
    0 \\ 1
  \end{pmatrix}
\end{equation}
hat die Koordinaten
\begin{equation}
  \begin{pmatrix}
    -\sin(\theta) \\ \cos(\theta)
  \end{pmatrix}
\end{equation}
in $\Sigma$.

Ein beliebiger Punkt $P$ mit Koordinaten in $\Sigma'$
\begin{equation}
  \begin{pmatrix}
    x' \\ y'
  \end{pmatrix} = x'\begin{pmatrix}
    1 \\ 0
  \end{pmatrix}+ y'\begin{pmatrix}
    0 \\ 1
  \end{pmatrix}
\end{equation}
hat in $\Sigma$ die folgenden Koordinaten
\begin{equation}
  \begin{split}
    \begin{pmatrix}
    x \\ y
  \end{pmatrix} &= x'\begin{pmatrix}
    \cos(\theta) \\ \sin(\theta)
  \end{pmatrix}+ y'\begin{pmatrix}
    -\sin(\theta) \\ \cos(\theta)
  \end{pmatrix} \\
  &=\begin{pmatrix}
    x'\cos(\theta)-y'\sin(\theta) \\ x'\sin(\theta)+y'\cos(\theta)
  \end{pmatrix}\\
  &= \begin{pmatrix}
    \cos(\theta) & -\sin(\theta) \\ 
    \sin(\theta) & \cos(\theta) 
  \end{pmatrix} \begin{pmatrix}
    x' \\ y'
  \end{pmatrix}. 
  \end{split}
\end{equation}
Mit $\mbf{r} \equiv \left( \begin{smallmatrix}
x \\ y
\end{smallmatrix}\right)$ in $\Sigma$ und $\mbf{r}'' \equiv \left( \begin{smallmatrix}
x' \\ y'
\end{smallmatrix}\right)$ in $\Sigma'$ erhalten wir also
\begin{equation}
  \eqbox{\mbf{r} = \mbf{R}(\theta )\mbf{r}',}
\end{equation}
wobei
\begin{equation}
  \eqbox{\mbf{R}(\theta ) \equiv \begin{pmatrix}
    \cos(\theta) & -\sin(\theta) \\ 
    \sin(\theta) & \cos(\theta) 
  \end{pmatrix}.}
\end{equation}
Die sogenannte \emph{Drehmatrix} oder \emph{Rotationsmatrix} $\mbf{R}(\theta )$ hat die folgenden Eigenschaften.
\begin{enumerate}
  \item Für $\theta=0$ ist die Rotationsmatrix die Einheitsmatrix:
  \begin{equation}
    \mbf{R}(0) \equiv \begin{pmatrix}
    \cos(0) & -\sin(0) \\ 
    \sin(0) & \cos(0) 
  \end{pmatrix} = \begin{pmatrix}
    1 & 0 \\ 
    0 & 1 
  \end{pmatrix} = \mathbb{I}.
  \end{equation}
  Dies entspricht der trivialen Transformation.
  \item Wird von $\Sigma$ nach $\Sigma'$ mit Drehmatrix $\mbf{R}(\alpha)$ gedreht und dann von $\Sigma'$ nach $\Sigma''$ mit $\mbf{R}(\beta)$, dann gilt
  \begin{equation}
    \begin{split}
      \mbf{r} &= \mbf{R}(\alpha)\mbf{r}' \\
      \mbf{r}' &= \mbf{R}(\beta)\mbf{r}''.
    \end{split}
  \end{equation}
  Also 
  \begin{equation}
      \mbf{r} = \mbf{R}(\alpha)\mbf{R}(\beta)\mbf{r}''.
  \end{equation}
  Hierbei berechnen wir
  \begin{equation}
    \begin{split}
      \mbf{R}(\alpha)\mbf{R}(\beta) &= \begin{pmatrix}
    \cos(\alpha) & -\sin(\alpha) \\ 
    \sin(\alpha) & \cos(\alpha) 
  \end{pmatrix}\begin{pmatrix}
    \cos(\beta) & -\sin(\beta) \\ 
    \sin(\beta) & \cos(\beta) 
  \end{pmatrix}\\
  &= \begin{pmatrix}
    \cos(\alpha)\cos(\beta)- \sin(\alpha)\sin(\beta)& -\cos(\alpha)\sin(\beta)-\sin(\alpha)\cos(\beta) \\ 
    \cos(\alpha)\sin(\beta)+\sin(\alpha)\cos(\beta) & - \sin(\alpha)\sin(\beta)+\cos(\alpha)\cos(\beta) 
  \end{pmatrix}\\
  &=\begin{pmatrix}
    \cos(\alpha+\beta) & -\sin(\alpha+\beta) \\ 
    \sin(\alpha+\beta) & \cos(\alpha+\beta) 
  \end{pmatrix}\\
  &= \mbf{R}(\alpha+\beta).
    \end{split}
  \end{equation}
  Also, $\mbf{r} = \mbf{R}(\alpha+\beta)\mbf{r}''$. 
  \item Falls die zweite Drehung um einen Winkel $\beta = -\alpha$ dreht, so gilt
  \begin{equation}
    \mbf{R}(\alpha)\mbf{R}(-\alpha) = \mbf{R}(0) = \mathbb{I}.
  \end{equation}
  Daher bildet $\mbf{R}(-\alpha)$ die inverse Matrix zu $\mbf{R}(\alpha)$:
  \begin{equation}
    \eqbox{\mbf{R}(\alpha)^{-1}=\mbf{R}(-\alpha).}
  \end{equation}
\end{enumerate}

Die Drehung in $2$ Dimensionen kann als eine Drehung um die $z$-Achse in $3$ Dimensionen interpretiert werden. Hier wird die $z$-Koordinate der Vektoren nicht verändert:
\begin{equation}
  \mbf{r} = \begin{pmatrix}
    \cos(\theta) & -\sin(\theta) & 0 \\ 
    \sin(\theta) & \cos(\theta) & 0 \\ 
    0 & 0 & 1
  \end{pmatrix}\mbf{r}'.
\end{equation}
\figbox
Wir bezeichnen die Drehmatrix, die diese Drehung darstellt als $\mbf{R}_z(\theta)$:
\begin{equation}
  \mbf{R}_z(\theta) = \begin{pmatrix}
    \cos(\theta) & -\sin(\theta) & 0 \\ 
    \sin(\theta) & \cos(\theta) & 0 \\ 
    0 & 0 & 1
  \end{pmatrix}.
\end{equation}
Drehungen um die $x-$ und $y-$ Achsen können auch konstruiert werden:
\begin{equation}
  \mbf{R}_x(\theta) = \begin{pmatrix}
    1&  0 & 0 \\
    0& \cos(\theta) & -\sin(\theta) \\ 
    0& \sin(\theta) & \cos(\theta)
  \end{pmatrix}
\end{equation}
und
\begin{equation}
  \mbf{R}_z(\theta) = \begin{pmatrix}
    \cos(\theta) & 0 &-\sin(\theta) \\ 
    0 & 1 & 0 \\
    \sin(\theta) & 0 & \cos(\theta) 
  \end{pmatrix}.
\end{equation}

Eine allgemeine Drehung in drei Dimensionen kann als Folge von drei Drehungen um verschiedene Achsen aufgefasst werden, zum Beispiel:
\begin{equation}
  \mbf{r} = \mbf{R}_z(\gamma)\mbf{R}_x(\beta)\mbf{R}_z(\alpha) \mbf{r}',
\end{equation}
wobei $\alpha$, $\beta$ und $\gamma$  als \emph{Eulerwinkel} bezeichnet werden.

Eine allgemeine, zeitunabhängige Transformation zwischen einem Bezugssystem $\Sigma$ und einem Bezugssystem $\Sigma'$, besteht aus einer Drehung und einem Versatz
\figbox
\begin{equation}
  \eqbox{\mbf{r} = \mbf{D}(\alpha,\beta,\gamma)\mbf{r}' + \mbf{r}_0.}
\end{equation}
Wie ändern sich die Bewegungsgleichungen? Wenn in $\Sigma'$
\begin{equation}
  m\ddot{\mbf{r}}' = \mbf{F}'
\end{equation}
 gilt, dann gilt in $\Sigma$
\begin{equation}
  m\ddot{\mbf{r}} = m\mbf{D}(\alpha,\beta,\gamma)\ddot{\mbf{r}}' = \mbf{D}(\alpha,\beta,\gamma)\mbf{F}' \equiv \mbf{F}
\end{equation}
mit dem Kraftfeld in $\Sigma$:
\begin{equation}
  \mbf{F} = \mbf{D}(\alpha,\beta,\gamma)\mbf{F}'.
\end{equation}
Ist $\Sigma'$ ein Inertialsystem, so ist auch $\Sigma$ ein Inertialsystem. (Kräftefreie Bewegung verläuft gleichförmig, das heisst, $\mbf{F}'=\mbf{0}\Rightarrow \mbf{F}=\mbf{0}$.)

Dieses Ergebnis gilt auch für gleichförmig gegeneinander bewegte Bezugssysteme: Sei
\begin{equation}\label{eq:gal}
  \eqbox{\mbf{r} = \mbf{D}\mbf{r}' + \mbf{v}_0 t + \mbf{r}_0}
\end{equation}
dann gilt
\begin{equation}
  m\ddot{\mbf{r}}' = \mbf{F}'
\end{equation}
und daher 
\begin{equation}
  m\ddot{\mbf{r}} = m \mbf{D}\ddot{\mbf{r}}' = \mbf{D}\mbf{F}' \equiv \mbf{F}.
\end{equation}

Eine Koordinatentransformation der Art (\ref{eq:gal}) wird als \emph{Galilei-Transformation} bezeichnet. Es gilt: Ist $\Sigma'$ ein Inertialsystem, so ist auch jedes durch eine Galilei-Transformation verbundenes Bezugssystem $\Sigma$ wieder ein Inertialsystem.

Für allgemeinere Transformationen gilt dies nicht mehr, das heisst, dass sich die Form der Bewegungsgleichungen ändert. Insbesondere können in beliebig bewegten Bezugssystemen \emph{Scheinkräfte} auftreten. Zwei wichtige Beispiele sind:
\begin{enumerate}
  \item \emph{Gleichförmig beschleunigte Bezugssysteme}
  \begin{equation}
    \begin{split}
      \mbf{r} &= \mbf{D}\mbf{r}' + \mbf{a}_0 \frac{t^2}{2}  + \mbf{v}_0 t + \mbf{r}_0 \\
      \dot{\mbf{r}} &= \mbf{D}\dot{\mbf{r}}' + \mbf{a}_0 t \\
      \ddot{\mbf{r}} &= \mbf{D}\ddot{\mbf{r}}' + \mbf{a}_0.
    \end{split}
  \end{equation}
  Falls in $\Sigma'$ die Gleichung $m\ddot{\mbf{r}}' = \mbf{F}'$ gilt, so gilt in $\Sigma$
  \begin{equation}
    m\ddot{\mbf{r}} = \mbf{D}\mbf{F}' + m\mbf{a}_0 \equiv \mbf{F} + \mbf{F}_a.
  \end{equation}
  In $\Sigma$ tritt $\mbf{F}_a = m\mbf{a}_0$ als eine zusätzliche Kraft auf, die proportional zur Masse $m$ ist.
  \item \emph{Gleichförmig rotierendes Bezugssystem}. In Inertialsystem $\Sigma$ rotiere ein Punkt $P$ mit Vektor $\mbf{e}$ (in $\Sigma$) um eine Drehachse $\mbf{n}$ (mit $\|\mbf{n}\| =1$) mit konstanter Winkelgeschwindigkeit $\omega$:
  \figbox
  Die Änderung von $\mbf{e}$ in einer kleinen Zeit $\Delta t$ wird $\Delta \mbf{e}$ genannt. Es gilt
  \begin{equation}
    \begin{split}
      \|\Delta \mbf{e}\| &= \|\mbf{e}\|\sin(\theta) \Delta \phi \\
      &= \|\mbf{e}\|\sin(\theta) \omega \Delta t. 
    \end{split}
  \end{equation}
  Der Vektor $\Delta \mbf{e}$ steht senkrecht zur $\mbf{e}$ und $\mbf{r}$:
  \begin{equation}
    \begin{split}
      \Delta \mbf{e} &\perp \mbf{e}\quad \text{und} \\
      \Delta \mbf{e} &\perp \mbf{n}.
    \end{split}
  \end{equation}
  Also 
  \begin{equation}
    \Delta\mbf{e} = (\mbf{n}\times \mbf{e}) \omega \Delta t.
  \end{equation}
  Formen wir diese Gleichung um, erhalten wir
  \begin{equation}
    \frac{\Delta \mbf{e}}{\Delta t} = \omega \mbf{n}\times \mbf{e}.
  \end{equation}
  Damit erhalten wir
  \begin{equation}
    \eqbox{\dot{\mbf{e}} = \boldsymbol{\omega}\times \mbf{e},}
  \end{equation}
  wobei 
  \begin{equation}
    \eqbox{\boldsymbol{\omega} \equiv \omega \mbf{n}}
  \end{equation}
  den \emph{Geschwindigkeitsvektoren} bildet.

Wir betrachten in einem Inertialsystem $\Sigma$ die Bahnkurve eines Teilchens  mit Positionsvektor
\begin{equation}
  \mbf{r}(t) = x(t) \hat{\mbf{x}} + y(t) \hat{\mbf{y}} +z(t) \hat{\mbf{z}}.
\end{equation}
In einem um eine Drehachse $\mbf{n}$ mit Winkelgeschwindigkeit $\omega$ rotierendes Bezugssystem $\Sigma'$ ist dieselbe Bahnkurve durch
\begin{equation}
  \mbf{r}(t) = x'(t) \mbf{e}_x'(t) + y'(t) \mbf{e}_y'(t) + z'(t) \mbf{e}_z'(t)
\end{equation}
gegeben, wobei $\mbf{e}_x'(t)$ der rotierende Einheitsvektor, der in die $x'$-Richtung bezüglich $\Sigma'$ zeigt, usw.

Wir berechnen den Geschwindigkeitsvektoren $\dot{\mbf{r}}$ wie folgt:
\begin{equation}
  \begin{split}
    \left(\frac{d\mbf{r}(t)}{dt}\right)_{\Sigma} &= \dot{x}'(t) \mbf{e}_x'(t) + \dot{y}'(t) \mbf{e}_y'(t) + \dot{z}'(t) \mbf{e}_z'(t) + x'(t) \dot{\mbf{e}}_x'(t) + y'(t) \dot{\mbf{e}}_y'(t) + z'(t) \dot{\mbf{e}}_z'(t) \\
    &= \left(\frac{d\mbf{r}(t)}{dt}\right)_{\Sigma'} + x'(t)(\boldsymbol{\omega}\times \mbf{e}_x') +z'(t)(\boldsymbol{\omega}\times \mbf{e}_z')+z'(t)(\boldsymbol{\omega}\times \mbf{e}_z') \\
    &= \left(\frac{d\mbf{r}(t)}{dt}\right)_{\Sigma'} + \boldsymbol{\omega}\times \mbf{r}.
  \end{split}
\end{equation}
Zusammengefasst: Der Geschwindigkeitsvektor bezüglich $\Sigma$ wird als die Summe von dem Geschwindigkeitsvektoren bezüglich $\Sigma'$ und einer Geschwindigkeit aufgrund der Rotation berechnet:
\begin{equation}
  \mbf{v}_{\Sigma} = \mbf{v}_{\Sigma'} + \boldsymbol{\omega}\times \mbf{r}.
\end{equation}
Damit folgt für die Beschleunigung in $\Sigma$
\begin{equation}
  \begin{split}
    \mbf{a}_{\Sigma} &\equiv \left(\frac{d\mbf{v}(t)}{dt}\right)_{\Sigma} = \left[\frac{d}{dt}(\mbf{v}_{\Sigma'} + \boldsymbol{\omega}\times \mbf{r})\right]_{\Sigma} \\
    &= \left(\frac{d\mbf{v}_{\Sigma'}(t)}{dt}\right)_{\Sigma'} + \boldsymbol{\omega}\times \left(\frac{d\mbf{r}(t)}{dt}\right)_{\Sigma} + \boldsymbol{\omega}\times \mbf{v}_{\Sigma'} + \boldsymbol{\omega}\times (\boldsymbol{\omega}\times \mbf{r}) \\
    &= \left(\frac{d\mbf{v}_{\Sigma'}(t)}{dt}\right)_{\Sigma'} + 2\boldsymbol{\omega}\times \mbf{v}_{\Sigma'}  + \boldsymbol{\omega}\times (\boldsymbol{\omega}\times \mbf{r})\\
    &= \mbf{a}_{\Sigma'} + 2\boldsymbol{\omega}\times \mbf{v}_{\Sigma'}  + \boldsymbol{\omega}\times (\boldsymbol{\omega}\times \mbf{r}).
  \end{split}
\end{equation}
Falls im Inertialsystem eine kräftefreie Bewegung verläuft
\begin{equation}
  m\mbf{a}_{\Sigma} = \mbf{0}
\end{equation}
so gilt in $\Sigma'$
\begin{equation}
  m(\mbf{a}_{\Sigma'} + 2\boldsymbol{\omega}\times \mbf{v}_{\Sigma'}  + \boldsymbol{\omega}\times (\boldsymbol{\omega}\times \mbf{r})) = \mbf{0}.
\end{equation}
Das heisst, die Bewegungsgleichung in $\Sigma'$ lautet
\begin{equation}
  \eqbox{m\ddot{\mbf{r}} = -2m\boldsymbol{\omega}\times \dot{\mbf{r}} - m\boldsymbol{\omega}\times (\boldsymbol{\omega}\times \mbf{r}).}
\end{equation}
\end{enumerate}

In rotierenden Bezugssystemen treten zwei Trägheitskräfte auf:
\begin{enumerate}
  \item \emph{Zentrifugalkraft}. Diese Kraft bezeichnet den Ausdruck
  \begin{equation}
    \mbf{F}_{\text{Z}} = - m\boldsymbol{\omega}\times (\boldsymbol{\omega}\times \mbf{r}).
  \end{equation}
  Die Zentrifugalkraft steht normal auf der Drehachse und zeigt radial nach aussen.
  \item \emph{Corioliskraft}. Diese Kraft beschreibt den Term
  \begin{equation}
    \mbf{F}_{\text{C}} = -2m\boldsymbol{\omega}\times \dot{\mbf{r}}.
  \end{equation}
  Die Corioliskraft steht normal auf der Drehachse und der Geschwindigkeit in $\Sigma'$.
\end{enumerate}